% =============================================================================
% BORRADOR – Partición de energía en impactos de nanocúmulos LJ
%
% Objetivo: física básica, sin motivación pedagógica.
% Audiencia: Physical Review E / Journal of Chemical Physics
%
% Compilar: pdflatex borrador_rebotes_gravedad.tex
% =============================================================================

\documentclass[12pt,a4paper]{article}

\usepackage[utf8]{inputenc}
\usepackage[T1]{fontenc}
\usepackage[spanish,es-nodecimaldot]{babel}
\usepackage{amsmath,amssymb}
\usepackage{graphicx}
\usepackage{hyperref}
\usepackage{siunitx}
\usepackage{booktabs}
\usepackage{physics}   % \dv, \pdv, \vb, \abs

% --- Macros físicas ----------------------------------------------------------
\newcommand{\cor}{e}
\newcommand{\coR}[1]{e_{#1}}
\newcommand{\ain}[1]{\alpha_{\mathrm{int},#1}}
\newcommand{\asu}[1]{\alpha_{\mathrm{sup},#1}}
\newcommand{\Ecm}[1]{E_{\mathrm{cm},#1}}
\newcommand{\Eint}[1]{E_{\mathrm{int},#1}}
\newcommand{\Tint}{T_{\mathrm{int}}}
\newcommand{\Tstar}{T^*}
\newcommand{\We}{\mathrm{We}}
\newcommand{\eps}{\varepsilon}
\newcommand{\epsps}{\varepsilon_{ps}}
\newcommand{\gstar}{g^*}
\newcommand{\lammps}{\textsc{Lammps}}
\newcommand{\TODO}[1]{\textbf{[TODO: #1]}}

\title{%
    \textbf{Partición de energía en el impacto de nanocúmulos\\
    Lennard-Jones sobre superficies FCC: cuatro canales,\\
    escalado con el tamaño del proyectil y adhesión}
}

\author{%
    % alternativa de correo: felipe.wasaff@uchile.cl
    Felipe Wasaff$^{1}$\thanks{felipe.wasaff.g@mail.pucv.cl}
    \quad y \quad
    Gonzalo Gutiérrez$^{2}$\\[4pt]
    \small $^{1}$Instituto de Matemáticas, Pontificia Universidad Católica
    de Valparaíso, Valparaíso, Chile\\
    \small $^{2}$Departamento de Física, Facultad de Ciencias,
    Universidad de Chile, Santiago, Chile
}

\date{\today}

% =============================================================================
\begin{document}
\maketitle

% =============================================================================
\begin{abstract}
% =============================================================================

Estudiamos mediante dinámica molecular (\lammps, unidades LJ) la partición de
energía en el primer rebote de un nanocúmulo de Lennard-Jones sobre una
superficie FCC cristalina bajo campo gravitatorio constante, reproduciendo y
extendiendo el trabajo de Gutiérrez et al.\ (2012). El sistema canónico es un
cúmulo FCC de $N=79$ átomos (5.ª capa completa, $r < 2.45\,\sigma$),
pre-minimizado, sobre una superficie de $8\times8\times6$ celdas unitarias
($a_0=1.55\,\sigma$) con interacción bola-superficie repulsiva ($c_{ps}=0$,
$T=0$). Identificamos cuatro canales energéticos cuya suma es exactamente
igual a la energía de impacto $E_0$:
\[
    \alpha_{\mathrm{cm}} + \alpha_{T_B} + \alpha_{\Delta V_B}
    + \alpha_{\mathrm{surf}} = 1,
\]
donde $\alpha_{T_B}\approx0.22$ es la fracción cinética interna de la bola
(constante para $v>0.9\,\sigma/\tau$), $\alpha_{\Delta V_B}\approx0.21$--$0.33$
es la variación neta de la PE intra-bola (expansión anarmónica post-impacto),
y $\alpha_{\mathrm{surf}}\approx0.07$--$0.16$ es la excitación directa de
fonones superficiales. Este desglose complementa el análisis original, que
caracterizó la energía colisional $U_c$ \emph{durante} el contacto: la
partición al primer apogeo separa por primera vez la expansión anarmónica
$\alpha_{\Delta V_B}$ de la excitación superficial, y muestra que la mayor
parte de la energía no recuperada permanece en la bola, no en la superficie.

Extendemos el estudio a $N\in\{1,2,7,13,55,79,87\}$ con interacción repulsiva.
La fracción interna total $\alpha_{\mathrm{bola}}=\alpha_{T_B}+\alpha_{\Delta V_B}$
sigue la ley de equipartición $\alpha_{\mathrm{bola}}=(3N-3)/(3N-3+3N_c)$ con
$N_c\approx N$ para $N\geq55$, prediciendo el límite asintótico
$\alpha_{\mathrm{bola}}\to\tfrac{1}{2}$ para $N\to\infty$.
Para $N=13$ (capa FCC completa más pequeña), la bola es demasiado rígida para
alcanzar el plateau de equipartición en el rango $v\leq2.0\,\sigma/\tau$ estudiado.

En simulaciones largas ($5000\,\tau$, $N=79$, $h_0=28\,\sigma$), la primera
altura de rebote aproxima bien la escala $h_1\approx T_B/(Mg)=8.3\,\sigma$
que aquí proponemos (ratio $=0.96$), pero el sistema no converge a ese
attractor: la altura CM decae más rápido que $T_B$, y tras 23 rebotes
($t\approx660\,\tau$) la bola queda atrapada sobre la superficie con
$h_n\lesssim1\,\sigma$.

Con interacción LJ completa ($c_{ps}=1$, adhesión), la bola pega en los
$4N\times6h_0=24$ casos estudiados ($v$ hasta $2.0\,\sigma/\tau$). La energía
de adhesión $E_{\rm adh}\propto N^{0.73}$ (área de contacto) supera
en todos los casos a la energía CM disponible tras el rebote; una
extrapolación de orden de magnitud sitúa el umbral en
$h_0^*(N)\sim260$--$440\,\sigma$ con $h_0^*\propto N^{-0.27}$.

\medskip
\noindent\textbf{Palabras clave:} dinámica molecular, nanocúmulos LJ,
coeficiente de restitución, partición energética, expansión anarmónica,
fonones de superficie, adhesión.
\end{abstract}

% =============================================================================
\section{Introducción}
\label{sec:intro}
% =============================================================================

La pelota que rebota es uno de los arquetipos más estudiados de la física
disipativa. A escala macroscópica, el coeficiente de restitución
$\cor = |v_f|/|v_i|$ es una propiedad del material y la energía disipada
va íntegramente a la superficie (o al medio). A escala nanométrica este
supuesto falla: un cúmulo de $N$ átomos tiene $3N-3$ grados de libertad
internos ---$3N-6$ vibracionales más 3 rotacionales--- que pueden absorber
energía durante cada impacto.

El trabajo de Gutiérrez et al.~\cite{Gutierrez2012} fue el primero en estudiar
esta \emph{cascada de disipación} mediante dinámica molecular: un cúmulo FCC
de argón de $N\approx79$ átomos bajo gravedad reducida $\gstar=0.026\,\sigma/\tau^2$,
con interacción bola-superficie puramente repulsiva. Sus resultados principales
fueron: (i)~la secuencia de alturas $h_n$ decrece más rápido que la serie
geométrica; (ii)~durante la colisión, la energía colisional bola-superficie
$U_c$ almacena transitoriamente la mayor parte de la energía ``disipada'',
más que la energía vibracional transferida a la superficie;
(iii)~para alturas de caída grandes la bola se funde durante el impacto
(un solo pico en $g(r)$ para $h_0>41.5\,\sigma$).
Sin embargo, su descomposición energética se evaluó \emph{durante} el
contacto (donde $U_c$ es un transitorio que se anula al separarse la bola):
no cuantificó la partición al primer apogeo, no separó la variación de la PE
intra-bola de la excitación directa de la superficie, y no estudió la
dependencia con $N$.

En este trabajo reproducimos ese experimento con \lammps{} usando un cúmulo de
$N=79$ átomos exactamente (capa FCC completa pre-minimizada, no creada con
\texttt{create\_atoms}) y extendemos el análisis en cuatro direcciones:
\textbf{(i)}~descomposición en \emph{cuatro} canales medidos de forma independiente;
\textbf{(ii)}~barrido $N\in\{1,2,7,13,55,79,87\}$ para calibrar la fórmula de
equipartición y determinar su escalado;
\textbf{(iii)}~simulación larga ($5000\,\tau$) para verificar la predicción
$h_1\approx T_B/(Mg)$ y el comportamiento a múltiples rebotes;
\textbf{(iv)}~barrido con adhesión ($c_{ps}=1$) para determinar el umbral
$h_0^*(N)$.

% =============================================================================
\section{Modelo y definiciones}
\label{sec:modelo}
% =============================================================================

\subsection{Hamiltoniano y descomposición energética}
\label{sec:hamiltoniano}

El sistema consta de un cúmulo de $N$ átomos (``proyectil'') y una superficie
FCC de $N_s$ átomos. El potencial LJ entre cualquier par a distancia $r$ es
\begin{equation}
    U_{\mathrm{LJ}}(r) = 4\eps \left[ \left(\frac{\sigma}{r}\right)^{12}
    - c_{ij}\left(\frac{\sigma}{r}\right)^{6} \right],
    \quad r < r_c = 2.5\,\sigma,
    \label{eq:lj}
\end{equation}
con $c_{ij}=1$ para pares bola-bola y superficie-superficie, y $c_{ij}=0$
(repulsivo puro, $c_{ps}=0$) o $c_{ij}=1$ (adhesión, $c_{ps}=1$) para
pares bola-superficie.

La energía interna de la bola se descompone en
\begin{equation}
    E_{\mathrm{int}} =
    \underbrace{\sum_{i=1}^{N}\tfrac{1}{2}m|\mathbf{v}_i-\mathbf{V}_{\rm cm}|^2}_{T_B}
    +
    \underbrace{\sum_{i<j}U_{\mathrm{LJ}}(r_{ij})}_{V_B}.
    \label{eq:TB_VB}
\end{equation}
Durante el vuelo libre $T_B$ y $V_B$ intercambian (oscilaciones internas)
mientras $E_{\rm int}=T_B+V_B$ permanece constante.

\subsection{Balance energético del primer rebote}
\label{sec:balance}

Denominamos $E_0 = Mg h_0^{\rm ef}$ a la energía mecánica inicial, donde
$h_0^{\rm ef}=Z_{\rm cm}(0)-z_0$ es la altura efectiva y $z_0$ es la posición
media de los mínimos de $Z_{\rm cm}$ (nivel de contacto). Midiendo cada
cantidad al primer apogeo (superíndice $(1)$) e inmediatamente antes del
impacto (referencia, superíndice $(0)$), definimos los cuatro canales:
\begin{align}
    \alpha_{\rm cm}         &= h_1/h_0^{\rm ef} = e_1^2, \label{eq:acm}\\
    \alpha_{T_B}            &= T_B^{(1)}/E_0,   \label{eq:atb}\\
    \alpha_{\rm surf}       &= \Delta(K_S + V_{\rm surf})/E_0, \label{eq:asurf}\\
    \alpha_{\Delta V_B}     &= 1-\alpha_{\rm cm}-\alpha_{T_B}-\alpha_{\rm surf}.
    \label{eq:avb}
\end{align}
La cuarta fracción $\alpha_{\Delta V_B}=(V_B^{(1)}-V_B^{(0)})/E_0$ puede
obtenerse de dos maneras: por sustracción (ec.~\eqref{eq:avb}, que garantiza
$\sum\alpha_i=1$ exactamente) o midiendo directamente $V_B$ en \lammps{}.
La diferencia entre ambas define el \emph{error de cierre}
$\varepsilon_{\rm cierre}=|\sum\alpha_i^{\rm medido}-1|$, que reportamos en
la Sección~\ref{sec:lammps} como control de calidad de la partición.

A partir del plateau observado $\alpha_{T_B}\approx0.22$ (Sección~\ref{sec:4canales}),
proponemos la escala de altura para el primer rebote
\begin{equation}
    h^* = T_B^{(1)}/(Mg),
    \label{eq:hstar}
\end{equation}
que predice $h_1\approx\alpha_{T_B}\,h_0^{\rm ef}$ cuando
$\alpha_{\rm cm}\approx\alpha_{T_B}$.

\subsection{Fórmula de equipartición}
\label{sec:equiparticion}

Proponemos que la fracción interna del primer rebote sigue una ley de
equipartición entre los grados de libertad activos durante el contacto:
\begin{equation}
    \alpha_{\rm int}(N) = \frac{3N-3}{3N-3+3N_c},
    \label{eq:equiparticion}
\end{equation}
donde $N_c$ es el número efectivo de osciladores de la huella de contacto.
El numerador cuenta los $3N-3$ grados de libertad internos de la bola
(vibracionales más rotacionales): un impacto no perfectamente central también
excita rotación, por lo que no restamos los 3 modos rotacionales. Para
$N\gtrsim13$ la diferencia entre $3N-3$ y $3N-6$ es menor que el resto de
las incertezas; para $N=2$ y $N=7$ no lo es, y esos tamaños se tratan solo
como controles cualitativos.
Si se identifica $\alpha_{\rm int}$ con $\alpha_{T_B}$ solo, se obtiene
$N_c\approx290$ (sin interpretación geométrica obvia). Si se identifica con
la fracción interna total $\alpha_{\rm bola}=\alpha_{T_B}+\alpha_{\Delta V_B}$
y se calibra con los datos de $N\geq55$, se obtiene $N_c\approx N$, lo que
predice el límite asintótico
\begin{equation}
    \alpha_{\rm bola}(N)\xrightarrow{N\to\infty}\frac{1}{2}.
    \label{eq:limite}
\end{equation}

% =============================================================================
\section{Detalles computacionales}
\label{sec:lammps}
% =============================================================================

Las simulaciones se realizan con \lammps~7~Feb~2024~\cite{Thompson2022}.

\paragraph{Superficie.}
Red FCC de $8\times8\times6$ celdas unitarias ($a_0=1.5496\,\sigma$),
condiciones de contorno periódicas en $xy$, no periódicas en $z$.
Las 2 capas inferiores están fijas (\texttt{setforce 0 0 0}).
Sin termostato (NVE global, $T=0$ inicial).

\paragraph{Cúmulo.}
El cúmulo de $N$ átomos se construye seleccionando las $N$ posiciones más
cercanas al origen de la red FCC del mismo parámetro de red. Las capas
esféricas completas corresponden a $N\in\{1,13,19,43,55,79,87,\ldots\}$;
de ellas simulamos $N\in\{1,13,55,79,87\}$, y agregamos $N=2$ y $N=7$
(capas incompletas) como controles del régimen molecular.
Se minimiza en vacío (algoritmo CG + FIRE) hasta fuerza máxima
$<10^{-6}\,\eps/\sigma$, obteniéndose energías de cohesión por átomo de
$-5.24\,\eps$ ($N=79$) y $-5.19\,\eps$ ($N=87$). Este protocolo garantiza
exactamente $N$ átomos para todas las alturas iniciales, a diferencia del
método \texttt{create\_atoms region sphere} de Gutiérrez et al.\ (2012),
que produce $N$ variable (79--84) dependiente de la posición del centro.
La bola se posiciona con velocidad CM nula a $h_0$ sobre el nivel de caja.

\paragraph{Gravedad y detector de rebotes.}
La gravedad $\gstar=0.026\,\sigma/\tau^2$ se aplica solo a los átomos de la
bola (\texttt{fix gravity}). $Z_{\rm cm}(t)$ se registra cada 20 pasos
($\Delta t_{\rm print}=0.1\,\tau$). Los mínimos y máximos de $Z_{\rm cm}$
se identifican en post-análisis con \texttt{scipy.signal.find\_peaks}.

\paragraph{Interacción bola-superficie.}
Para $c_{ps}=0$: tabla LJ repulsiva pre-generada ($4/r^{12}$ puro, con
fuerza analítica $F=-dV/dr$), tabulada en $r\in[0.75,2.5]\,\sigma$ con
5000 puntos e interpolación lineal (\texttt{pair\_style table}).
Para $c_{ps}=1$: LJ completo con \texttt{pair\_coeff 1 2 lj/cut 1.0 1.0 2.5}.
El cambio se controla con la variable \lammps{} \texttt{c\_ps} en
\texttt{in.bounce\_cluster}.

\paragraph{Integración y conservación de energía.}
Las ecuaciones de movimiento se integran con velocity-Verlet
(\texttt{fix nve}) y paso de tiempo $\Delta t = 5\times10^{-3}\,\tau$
(el estudio original~\cite{Gutierrez2012} usó Beeman con
$\Delta t = 4.65\times10^{-4}\,\tau$). Todos los potenciales se truncan
\emph{y desplazan} en $r_c=2.5\,\sigma$ (\texttt{pair\_modify shift yes});
el desplazamiento no altera fuerzas ni trayectorias, pero es esencial para
el balance energético: sin él, cada par que cruza $r_c$ introduce un salto
de $\pm|U(r_c)|=0.016\,\eps$ en la energía potencial registrada, y el error
acumulado alcanza 2--6\% de $E_0$ por impacto. Con potenciales desplazados,
la energía total (incluida la PE gravitatoria) se conserva a
$|\Delta E|<0.03\,\eps$ ($\sim4\times10^{-4}\,E_0$) en $120\,\tau$
incluyendo un impacto completo, lo que valida $\Delta t=5\times10^{-3}\,\tau$.
Como control de la partición, medimos $\alpha_{\Delta V_B}$ de forma
independiente (directamente de $V_B$) además de por sustracción
(ec.~\eqref{eq:avb}): el error de cierre resultante es
$|\sum\alpha_i-1|<10^{-4}$ para el primer rebote de $N=79$,
$h_0=28.15\,\sigma$.

% =============================================================================
\section{Resultados}
\label{sec:resultados}
% =============================================================================

% -----------------------------------------------------------------------------
\subsection{Validación cruzada: reproducción de Gutiérrez et al.\ (2012)}
\label{sec:reproduccion}
% -----------------------------------------------------------------------------

Reproducimos los parámetros originales ($c_{ps}=0$, $T=0$, $\gstar=0.026$,
6 alturas desde $h_0=11.03$ hasta $79.53\,\sigma$) con nuestro cúmulo
pre-minimizado de $N=79$ átomos. El nivel de contacto efectivo es
$z_0\approx11.5\,\sigma$ (altura del CM en los mínimos).

El COR del primer rebote va de $e_1=0.92$ ($h_0^{\rm ef}=8.8\,\sigma$) a
$e_1=0.59$ ($h_0^{\rm ef}=79.0\,\sigma$). La referencia~\cite{Gutierrez2012}
no reporta valores numéricos de $e_1$, pero sí las secuencias de alturas
$h_n$ (su Fig.~7), que digitalizamos para la comparación cuantitativa de la
Tabla~\ref{tab:comparacion2012}. La digitalización revela además una
diferencia de convención: en~\cite{Gutierrez2012} el parámetro $h_0$ es la
altura \emph{absoluta} del CM (las curvas de su Fig.~7 parten exactamente
en $h_0$ y convergen al nivel de reposo $z_{\rm rest}\approx10.3\,\sigma$),
mientras que aquí $h_0$ se mide desde el tope de la caja de la superficie.
Por eso la comparación debe hacerse a altura efectiva de caída
$h_0^{\rm ef}$ comparable, no a $h_0$ nominal igual.

\begin{table}[htbp]
\centering
\small
\caption{Comparación con la Fig.~7 de~\cite{Gutierrez2012} (valores
digitalizados; incerteza $\approx\pm0.7\,\sigma$ en las alturas, que domina
el error de $e_1$). Las alturas efectivas de 2012 son
$h_0^{\rm ef}=h_0-z_{\rm rest}$ con $z_{\rm rest}\approx10.3\,\sigma$.
La fila $h_0=11.03$ de 2012 es una caída de $<1\,\sigma$ (cuasi-elástica
trivial) y no admite un $e_1$ confiable.}
\label{tab:comparacion2012}
\begin{tabular}{c|ccc|cc}
\toprule
 & \multicolumn{3}{c|}{Gutiérrez et al.\ 2012 (Fig.~7)} &
 \multicolumn{2}{c}{Este trabajo ($N=79$)} \\
$h_0$ ($\sigma$) & $h_0^{\rm ef}$ & $h_1^{\rm ef}$ & $e_1$ &
$h_0^{\rm ef}$ & $e_1$ \\
\midrule
11.03 &  0.6 &  0.4 & ---  &  8.81 & 0.924 \\
18.64 &  8.4 &  6.8 & 0.90 & 17.21 & 0.708 \\
28.15 & 17.9 &  3.3 & 0.43 & 25.91 & 0.552 \\
41.47 & 31.3 &  8.2 & 0.51 & 41.80 & 0.633 \\
58.60 & 48.4 & 14.3 & 0.54 & 56.93 & 0.614 \\
79.53 & 69.2 & 18.2 & 0.51 & 79.03 & 0.587 \\
\bottomrule
\end{tabular}
\end{table}

Comparando a $h_0^{\rm ef}$ comparable, el acuerdo es excelente en el
régimen cuasi-elástico ($e_1=0.90$ en 2012 con $h_0^{\rm ef}=8.4\,\sigma$
frente a $0.92$ aquí con $8.8\,\sigma$) y bueno a alta energía, donde los
valores de 2012 quedan sistemáticamente $\sim0.06$--$0.09$ por debajo
(p.~ej.\ $0.51$ con $h_0^{\rm ef}=69\,\sigma$ frente a $0.59$ con
$79\,\sigma$). Esa diferencia es atribuible a que el cúmulo de 2012 se
generó con \texttt{create\_atoms} sin pre-minimización ($N$ variable
79--84 y energía interna residual, que abre canales de absorción
adicionales) y a que el integrador difiere (Beeman en 2012,
velocity-Verlet aquí), por lo que solo cabe esperar acuerdo estadístico.
El caso $h_0^{\rm ef}=17.9\,\sigma$ de 2012 ($e_1=0.43$) es un outlier
frente a la tendencia de ambos conjuntos; nótese que nuestra serie también
presenta no monotonía en esa zona (mínimo $e_1=0.55$ en
$h_0^{\rm ef}=25.9\,\sigma$, Tabla~\ref{tab:particion4}), lo que sugiere
sensibilidad genuina a la fase interna de la bola al impactar.
La caída más rápida que la predicción geométrica $e_1=\mathrm{cte}$
confirma la cascada de disipación en ambos conjuntos de datos.

% -----------------------------------------------------------------------------
\subsection{Partición en cuatro canales para \texorpdfstring{$N=79$}{N=79}}
\label{sec:4canales}
% -----------------------------------------------------------------------------

La Tabla~\ref{tab:particion4} muestra la descomposición del primer rebote para
$N=79$, $c_{ps}=0$, $T=0$. El resultado central es que para
$v_{\rm imp}\geq0.9\,\sigma/\tau$, la expansión anarmónica supera ampliamente
a la excitación superficial directa: $\alpha_{\Delta V_B}\approx21$--$33\%$
frente a $\alpha_{\rm surf}\approx7$--$16\%$.

\begin{table}[ht]
\centering
\small
\caption{Partición de la energía del primer rebote en cuatro canales para
$N=79$ (cúmulo FCC pre-minimizado), $c_{ps}=0$, $T=0$.
$\sum\alpha_i=1$ por construcción.}
\label{tab:particion4}
\begin{tabular}{ccccccc}
\toprule
$h_0^{\rm ef}$ & $v_{\rm imp}$ & $e_1$ &
$\alpha_{\rm cm}$ & $\alpha_{T_B}$ & $\alpha_{\Delta V_B}$ &
$\alpha_{\rm surf}$ \\
($\sigma$) & ($\sigma/\tau$) & & & & & \\
\midrule
 8.81 & 0.677 & 0.924 & 0.853 & 0.021 & 0.056 & 0.070 \\
17.21 & 0.946 & 0.708 & 0.502 & 0.213 & 0.208 & 0.077 \\
25.91 & 1.161 & 0.552 & 0.304 & 0.319 & 0.310 & 0.066 \\
41.80 & 1.474 & 0.633 & 0.401 & 0.257 & 0.264 & 0.078 \\
56.93 & 1.721 & 0.614 & 0.378 & 0.197 & 0.328 & 0.098 \\
79.03 & 2.027 & 0.587 & 0.345 & 0.204 & 0.294 & 0.156 \\
\bottomrule
\end{tabular}
\end{table}

La fracción cinética interna $\alpha_{T_B}$ es prácticamente constante para
$v_{\rm imp}\geq0.9\,\sigma/\tau$ ($\alpha_{T_B}\approx0.21$--$0.24$).
Este plateau cuantifica, al primer apogeo, el calentamiento interno de la
bola que~\cite{Gutierrez2012} describió cualitativamente durante el contacto.
La fracción interna total $\alpha_{\rm bola}=\alpha_{T_B}+\alpha_{\Delta V_B}$
varía entre 0.42 y 0.63, superando en todos los casos a $\alpha_{\rm cm}$
para $v>1.0\,\sigma/\tau$.

La excitación superficial directa $\alpha_{\rm surf}$ es la contribución
más pequeña del sistema, entre 7 y 16\% de $E_0$.
Este resultado matiza la imagen que sugiere el análisis durante el contacto
de~\cite{Gutierrez2012}: allí la energía colisional $U_c$ domina el balance
\emph{transitorio}, pero $U_c$ se anula cuando la bola se separa. La
partición al apogeo muestra que la mayor parte de esa energía regresa a la
bola como expansión anarmónica ($\alpha_{\Delta V_B}$) y no queda en la
superficie. Ambas descripciones son compatibles: miden la distribución en
instantes distintos del rebote.

La Figura~\ref{fig:secuencia} ilustra esta secuencia con instantáneas de la
simulación ($h_0=28.15\,\sigma$): en la compresión máxima la energía está
mayoritariamente almacenada como PE intra-bola y de contacto (transitorias),
mientras que al primer apogeo la bola retiene $\sim$62\% de $E_0$ como
energía interna ($T_B+\Delta V_B$) ---visible como el ``moteado'' caliente
de sus átomos--- y la superficie solo $\sim$6\%.

\begin{figure}[htbp]
    \centering
    \includegraphics[width=\textwidth]{../../../data/outputs/terminado/fig_sequence.png}
    \caption{Cuatro instantes del primer rebote ($N=79$, $h_0=28.15\,\sigma$,
    $c_{ps}=0$): \textbf{(a)}~caída libre, \textbf{(b)}~compresión máxima,
    \textbf{(c)}~expansión post-impacto y \textbf{(d)}~primer apogeo.
    Los átomos de la bola se colorean por su KE interna (sin el movimiento
    del CM, escala $0$--$0.30\,\eps$) y los de la superficie por su KE
    ($0$--$0.01\,\eps$, escala $\times30$ más sensible para visualizar los
    fonones). La barra bajo cada panel muestra la partición instantánea de
    $E_0$; $U_c$ (gris) solo es apreciable durante el contacto, y en el
    apogeo la barra reproduce los cuatro canales de la
    Tabla~\ref{tab:particion4} (con $U_F=\alpha_{\rm cm}E_0$).}
    \label{fig:secuencia}
\end{figure}

% -----------------------------------------------------------------------------
\subsection{Escalado con \texorpdfstring{$N$}{N}: calibración de \texorpdfstring{$N_c$}{Nc}}
\label{sec:sweep_N}
% -----------------------------------------------------------------------------

Corremos las mismas seis alturas con $N\in\{1,2,7,13,55,79,87\}$, usando
cúmulos pre-minimizados en capas FCC completas y promediando sobre los puntos
con $v_{\rm imp}>1.1\,\sigma/\tau$. Los resultados se muestran en la
Tabla~\ref{tab:sweep_N} y la Figura~\ref{fig:sweep_N}.

\begin{table}[ht]
\centering
\caption{Fracciones del primer rebote promediadas sobre $v_{\rm imp}>1.1\,\sigma/\tau$,
en función de $N$. Las dos últimas columnas muestran la predicción de la
ec.~\eqref{eq:equiparticion} calibrada con $\alpha_{T_B}$ ($N_c=290$) y con
$\alpha_{\rm bola}$ ($N_c=N$).}
\label{tab:sweep_N}
\begin{tabular}{cccccccc}
\toprule
$N$ & capas & $\langle e_1\rangle$ & $\alpha_{T_B}$ & $\alpha_{\rm bola}$ &
$\alpha_{\rm surf}$ &
$\hat\alpha_{T_B}(N_c{=}290)$ & $\hat\alpha_{\rm bola}(N_c{=}N)$ \\
\midrule
 1 & 1 (completa) & — & 0.000 & 0.000 & —     & 0.000 & 0.000 \\
 2 & 1+1          & 0.902 & 0.002 & 0.003 & 0.183 & 0.003 & 0.333 \\
 7 & 1+6/12$^a$   & 0.840 & 0.034 & 0.093 & 0.198 & 0.020 & 0.462 \\
13$^b$ & 2 (completa) & 0.735 & 0.066 & 0.107 & 0.339 & 0.040 & 0.480 \\
55 & 4 (completa) & 0.521 & 0.217 & 0.474 & 0.254 & 0.157 & 0.495 \\
79 & 5 (completa) & 0.597 & 0.244 & 0.543 & 0.100 & 0.212 & 0.497 \\
87 & 6 (completa) & 0.594 & 0.250 & 0.543 & 0.103 & 0.229 & 0.497 \\
\bottomrule
\end{tabular}
\footnotesize $^a$ N=7 es una capa incompleta (7 de 13 posibles en la 2.ª capa FCC);
sus valores presentan mayor dispersión.\\
\footnotesize $^b$ N=13 no alcanza el plateau en $v\leq2.0\,\sigma/\tau$
(ver texto); su promedio no es representativo y se excluye de la
calibración de $N_c$.\\
\footnotesize Para $N=1$ no existen modos internos ($\alpha_{T_B}\equiv0$);
la fila es un control. Para $N=2$, casi toda la energía no recuperada va a
la superficie ($\alpha_{\rm surf}=0.183$, con
$\alpha_{\rm cm}=\langle e_1\rangle^2\approx0.81$): el dímero apenas se
acopla a sus 3 modos internos en impactos verticales.
\end{table}

\paragraph{Calibración de $N_c$.}
Para $\alpha_{T_B}$ solo (ec.~\eqref{eq:equiparticion}), el ajuste da
$N_c^{\rm ef}\in\{169,195,241,258\}$ para $N\in\{13,55,79,87\}$,
creciendo débilmente con $N$ ($N_c\propto N^{0.20}$) sin una interpretación
geométrica directa. En cambio, usando $\alpha_{\rm bola}$ en
ec.~\eqref{eq:equiparticion} se obtiene $N_c\approx N$ para $N\geq55$
(columna $\hat\alpha_{\rm bola}$). El acuerdo no es perfecto: los datos
($\alpha_{\rm bola}=0.543$ para $N=79$ y $87$) superan sistemáticamente
la predicción con $N_c=N$ (0.497) e incluso la cota asintótica $1/2$,
lo que en sentido estricto implica $N_c<N$: la inversión de la
ec.~\eqref{eq:equiparticion} da $N_c/N=1.09$, $0.83$ y $0.83$ para
$N=55$, $79$ y $87$. El exceso sistemático puede originarse en
la energía rotacional de la bola (incluida en $\alpha_{\rm bola}$ pero
no en el conteo de osciladores armónicos) o en una subestimación de
$\alpha_{\rm surf}$; discriminar ambas requiere medir el momento angular
post-rebote (\texttt{compute angmom}), pendiente para una revisión.
Con esa salvedad, la identificación $N_c\sim N$ ---cada átomo de la bola
acoplado efectivamente a un oscilador de la superficie--- captura la
tendencia y predice el límite $\alpha_{\rm bola}\to1/2$ para $N\to\infty$.

\paragraph{Régimen de bola pequeña ($N=13$).}
El cúmulo de $N=13$ (única capa esférica FCC completa de radio $1.10\,\sigma$)
no alcanza el plateau de equipartición en el rango $v\leq2.0\,\sigma/\tau$:
$\alpha_{T_B}$ sigue creciendo desde $0.003$ (baja energía) hasta $0.152$
(máxima velocidad estudiada), mientras que la predicción de equipartición
daría $\alpha_{T_B}\approx0.04$ ($N_c=290$) o $\alpha_{\rm bola}\approx0.48$
($N_c=N$). La rigidez del icosaedro de una sola capa impide la termalización
rápida que permite el plateau en cúmulos más grandes.

\begin{figure}[htbp]
    \centering
    \includegraphics[width=\textwidth]{../../../data/outputs/terminado/fig_alpha_TB_full_sweep.pdf}
    \caption{Barrido $N\in\{1,2,7,13,55,79,87\}$, $c_{ps}=0$, $T=0$.
    \textbf{(a)}~$\alpha_{T_B}$ vs $v_{\rm imp}$ para cada $N$; la línea
    vertical punteada marca $v^*=1.1\,\sigma/\tau$ (inicio del plateau).
    \textbf{(b)}~$\langle\alpha_{T_B}\rangle_{\rm plateau}$ vs $N$ (las barras
    de error son la desviación estándar sobre los puntos con $v>v^*$),
    comparado con la fórmula ec.~\eqref{eq:equiparticion} para $N_c=290$
    (curva azul sólida) y con el ajuste $N_c(N)=99\,N^{0.20}$ (curva negra
    discontinua). El símbolo abierto de $N=13$ indica que ese tamaño no
    alcanza el plateau. El inset muestra $N_c^{\rm ef}(N)$ y su ajuste de
    ley de potencias.}
    \label{fig:sweep_N}
\end{figure}

% -----------------------------------------------------------------------------
\subsection{Simulación larga: verificación de la predicción \texorpdfstring{$h^*=T_B/(Mg)$}{h*}}
\label{sec:largo}
% -----------------------------------------------------------------------------

Para poner a prueba la escala $h^*=T_B/(Mg)$ propuesta en la
Sección~\ref{sec:balance} ---y en particular si la altura de rebote
\emph{converge} a ella---, corremos una simulación de $5000\,\tau$
($N=79$, $h_0=28.15\,\sigma$, $c_{ps}=0$, mismos parámetros).

\paragraph{Primer rebote.}
La altura medida $h_1=7.91\,\sigma$ está en excelente acuerdo con la predicción
$h^*(T_B^{(1)})=T_B^{(1)}/(Mg)=8.27\,\sigma$ (ratio $=0.96$).
Este resultado confirma que la relación $h_1\approx T_B/(Mg)$ es válida para
el primer impacto.

\paragraph{Rebotes sucesivos y colapso.}
A diferencia de lo que sugiere la nomenclatura «$h^*$» como attractor
estacionario, el ratio $h_n/h^*(T_B^{(n)})$ decrece monotónamente rebote a
rebote (de 0.96 en $n=1$ a $\sim0.18$ en $n=20$--$23$).
La altura CM $h_n$ decae geométricamente ($e_1^{\rm ef}\approx0.97$ por rebote),
mientras que $T_B$ disminuye más despacio (de $17.0\,\eps$ en $n=1$ a
$11.8\,\eps$ en $n=23$, factor $\times0.7$).

Tras 23 rebotes ($t\approx660\,\tau$), la altura de apogeo se ha reducido a
$h_n\approx1.0\,\sigma$: la bola ya no se separa efectivamente de la
superficie. En $t\approx713\,\tau$ el CM cae por debajo de $z_0$
(colapso): átomos de la bola se filtran entre los huecos del lattice de
la superficie al no haber barrera repulsiva superior en la pared $+z$.

\paragraph{Interpretación.}
La predicción $h_1\approx h^*=T_B/(Mg)$ es una consecuencia del primer impacto
($\alpha_{\rm cm}\approx\alpha_{T_B}$ en el primer rebote, ver
Tabla~\ref{tab:particion4}), no un attractor dinámico. En rebotes sucesivos,
la energía CM adicional disponible (caída desde $h_n\ll h_0$) es pequeña,
por lo que $T_B$ crece poco mientras $h_n$ sigue cayendo por inelasticidad.

\begin{figure}[htbp]
    \centering
    \includegraphics[width=\textwidth]{../../../data/outputs/terminado/fig_paso_B_convergencia.pdf}
    \caption{Simulación larga ($5000\,\tau$, $N=79$, $h_0=28.15\,\sigma$,
    $c_{ps}=0$).
    \textbf{(a)}~Trayectoria $Z_{\rm cm}(t)-z_0$ durante los primeros
    $1000\,\tau$; la línea roja discontinua marca el colapso a través de la
    red en $t\approx713\,\tau$ (artefacto del potencial repulsivo puro).
    El inset muestra los $5000\,\tau$ completos.
    \textbf{(b)}~Alturas de apogeo $h_n$ (círculos azules) y predicción
    $h^*_n=T_B^{(n)}/(Mg)$ (cuadrados rojos) por número de rebote.
    \textbf{(c)}~Razón $h_n/h^*_n$: el valor del primer rebote (0.96) es el
    único que se aproxima a la unidad (línea discontinua).}
    \label{fig:largo}
\end{figure}

% -----------------------------------------------------------------------------
\subsection{Adhesión: barrido con \texorpdfstring{$c_{ps}=1$}{c\_ps=1}}
\label{sec:adhesion}
% -----------------------------------------------------------------------------

Con interacción bola-superficie LJ completa ($c_{ps}=1$, $\eps_{ps}=\eps$),
corremos los mismos $N\in\{13,55,79,87\}$ a las 6 alturas
($v_{\rm imp}=0.71$--$2.03\,\sigma/\tau$). La bola pega en los 24/24 casos
(Tabla~\ref{tab:adhesion}): no se detecta ningún apogeo post-contacto, y
$Z_{\rm cm}$ al final de la simulación ($400\,\tau$) se sitúa a
$\approx10\,\sigma$, sobre el tope de la superficie.

\begin{table}[ht]
\centering
\caption{Energía de adhesión y umbral estimado para $c_{ps}=1$,
$\eps_{ps}=\eps$.
$E_{\rm adh}=|\langle U_c\rangle_{\rm reposo}|$ (promedio sobre los últimos
50 pasos de cada simulación). $h_0^*$ estimado del balance energético
$\alpha_{\rm cm}\,E_0^*=E_{\rm adh}$ con $\alpha_{\rm cm}\approx0.35$.}
\label{tab:adhesion}
\begin{tabular}{ccccc}
\toprule
$N$ & $E_{\rm adh}$ ($\eps$) & $E_{\rm adh}/N$ ($\eps$/át.) &
$h_0^*$ estimado ($\sigma$) & Escala \\
\midrule
13 & 52  & 4.00 & 439 &  \\
55 & 135 & 2.45 & 270 &  \\
79 & 197 & 2.49 & 274 & $h_0^*\propto N^{-0.27}$ \\
87 & 209 & 2.40 & 264 &  \\
\bottomrule
\end{tabular}
\end{table}

\paragraph{Energía de adhesión.}
$E_{\rm adh}$ escala como $N^{0.73}$ (Figura~\ref{fig:adhesion}, panel
derecho), próximo al escalado del área de contacto $\propto N^{2/3}=N^{0.67}$.
La energía por átomo $E_{\rm adh}/N\approx2.4$--$4.0\,\eps$ decae con $N$,
reflejando que los átomos del interior de la bola contribuyen menos al
contacto que los átomos superficiales.

\paragraph{Umbral de rebote.}
El balance energético para el umbral requiere que la energía CM disponible
tras el rebote iguale la energía de adhesión:
\begin{equation}
    h_0^*(N) = \frac{E_{\rm adh}(N)}{\alpha_{\rm cm}\,Mg}
             = \frac{E_{\rm adh}(N)}{0.35\,N\gstar}.
    \label{eq:h0star}
\end{equation}
Los valores estimados ($h_0^*\approx260$--$440\,\sigma$, Tabla~\ref{tab:adhesion})
están fuera del rango de alturas estudiado ($h_{0,\rm max}=79.5\,\sigma$).
La dependencia $h_0^*\propto N^{-0.27}$ indica que cúmulos \emph{más grandes}
tienen umbrales \emph{más bajos}: su masa crece como $N$ pero la adhesión
solo como $N^{0.73}$, haciendo más fácil el escape para proyectiles grandes.

Este umbral debe leerse como cota de \emph{orden de magnitud}, no como
predicción cuantitativa, por tres razones. Primero, es una extrapolación a
$3$--$5$ veces la altura máxima validada. Segundo, a esas alturas
($v_{\rm imp}\sim3.7$--$4.8\,\sigma/\tau$) el cúmulo ya no rebota como
sólido: \cite{Gutierrez2012} observó fusión de la bola para
$h_0>41.5\,\sigma$, de modo que en el régimen del umbral el área de
contacto ---y con ella $E_{\rm adh}$--- crecería con $v$, en vez de ser
constante como supone la ec.~\eqref{eq:h0star}. Tercero,
$\alpha_{\rm cm}\approx0.35$ está tomado del caso repulsivo; con adhesión
la retención de energía CM será menor. Los tres sesgos apuntan en la misma
dirección (retrasar el rebote), por lo que $h_0^*$ es probablemente una
cota inferior. Verificarlo requiere simulaciones con $c_{ps}=1$ a
$h_0$ intermedias-altas midiendo $E_{\rm adh}(v)$.

\begin{figure}[htbp]
    \centering
    \includegraphics[width=\textwidth]{../../../data/outputs/terminado/fig_paso_C_adhesion.pdf}
    \caption{\textbf{(a)}~Trayectorias $Z_{\rm cm}(t)$ para $N=79$,
    $c_{ps}=1$, 6 alturas (primeros $120\,\tau$; no hay más eventos después).
    Todas terminan en $Z_{\rm cm}\approx10\,\sigma$ (tope de la superficie,
    línea gris discontinua): adhesión en los 6 casos.
    \textbf{(b)}~Umbral $h_0^*(N)$ estimado de la ec.~\eqref{eq:h0star}
    (círculos rojos; la barra de error propaga la dispersión de
    $E_{\rm adh}$ entre alturas), comparado con el ajuste
    $E_{\rm adh}=7.9\,N^{0.73}$ (discontinua negra) y con la escala
    $\propto N^{-1/3}$ (punteada azul).}
    \label{fig:adhesion}
\end{figure}

% =============================================================================
\section{Discusión}
\label{sec:discusion}
% =============================================================================

\subsection{Origen de \texorpdfstring{$\alpha_{\Delta V_B}$}{α\_ΔVB}:
expansión anarmónica post-impacto}
\label{sec:anarmonico}

El canal $\alpha_{\Delta V_B}$ no tiene análogo en modelos de esferas rígidas.
Durante la compresión, el impulso de la superficie convierte KE del CM en PE
intra-bola. Al rebotar, el CM recupera KE, pero la relajación de la PE interna
hacia el equilibrio tarda varias vibraciones ($\sim1$--$5\,\tau$), tiempo en que
el cúmulo ya se aleja. En el primer apogeo, $V_B>V_B^{(0)}$, y la diferencia
registra la energía potencial que la bola «se lleva».

Esta asimetría tiene dos consecuencias: (i)~cualquier partición que use solo
la KE interna $T_B^{(1)}$ subestima en un factor $\sim2$ el calor depositado
en la bola; (ii)~la cantidad natural para el balance de equipartición es
$\alpha_{\rm bola}=\alpha_{T_B}+\alpha_{\Delta V_B}$, no $\alpha_{T_B}$.

\subsection{Nc efectivo y escalado con N}
\label{sec:Nc_interpretacion}

La calibración $N_c\approx N$ a partir de $\alpha_{\rm bola}$ para $N\geq55$
tiene una lectura directa: la energía de impacto se distribuye entre $3N-3$
modos internos de la bola y $3N$ modos superficiales activos, es decir, la
huella activa del impacto involucra efectivamente $N$ átomos de la superficie.
Esto es consistente con que el radio de la bola ($\sim N^{1/3}$) determine
también la profundidad de penetración de la onda de presión en la superficie,
activando $\sim N$ átomos en volumen (no solo en área de contacto $\sim N^{2/3}$).

Para $N=13$, el sistema no alcanza el régimen de equipartición: la bola es
demasiado rígida (una sola capa esférica) y $\alpha_{\rm bola}\approx0.11$
queda muy por debajo de la predicción de $\sim0.48$. El resultado marca una
transición entre el régimen ``molecular'' ($N\leq13$) y el régimen ``bulk-like''
($N\geq55$).

\subsection{Dinámica a largo plazo y colapso}
\label{sec:largo_discusion}

La simulación de $5000\,\tau$ revela que $h^*=T_B/(Mg)$ no es un attractor
estacionario sino una cota aproximada válida solo para el primer rebote
(cuando $\alpha_{\rm cm}\approx\alpha_{T_B}$). A medida que $h_n$ decrece
geométricamente y $T_B$ crece más lentamente, la bola se aproxima a la
superficie hasta quedar atrapada sobre ella. El colapso a través del lattice
de la superficie (detectable por $Z_{\rm cm}<z_0$) es un artefacto del
potencial repulsivo puro que no tiene barrera lateral: con $c_{ps}=1$
este colapso no ocurre porque la bola queda adherida antes de esa etapa.

\subsection{Umbral de adhesión y soft landing}
\label{sec:adhesion_discusion}

Los resultados de $c_{ps}=1$ muestran que la energía de adhesión LJ es
considerablemente mayor que la energía CM disponible para todos los $N$
y $h_0$ estudiados. El umbral $h_0^*(N)\propto N^{-0.27}$ tiene una implicación
directa para aplicaciones de \emph{soft landing}: cúmulos más grandes requieren
menores velocidades de impacto relativas para adherirse ($v^*\propto N^{-0.13}$),
lo que favorece la deposición de cúmulos grandes a velocidades intermedias.

% =============================================================================
\section{Conclusiones}
\label{sec:conclusiones}
% =============================================================================

Hemos reproducido y extendido el estudio de Gutiérrez et al.~\cite{Gutierrez2012}
con \lammps{} usando cúmulos FCC pre-minimizados de tamaño exacto
($N\in\{1,2,7,13,55,79,87\}$). Las conclusiones principales son:

\begin{enumerate}

    \item \textbf{El impacto involucra cuatro canales energéticos distintos.}
    $\alpha_{\rm cm}+\alpha_{T_B}+\alpha_{\Delta V_B}+\alpha_{\rm surf}=1$,
    con $\alpha_{T_B}\approx0.22$ constante, $\alpha_{\Delta V_B}\approx0.21$--$0.33$
    y $\alpha_{\rm surf}\approx0.07$--$0.16$ para $N=79$, $v>0.9\,\sigma/\tau$.
    La expansión anarmónica domina sobre la excitación superficial directa en
    factor $\sim3$--$4$.

    \item \textbf{La fracción interna total $\alpha_{\rm bola}$ escala según
    equipartición con $N_c\approx N$ para $N\geq55$.}
    Esto predice $\alpha_{\rm bola}\to1/2$ para $N\to\infty$, aunque los
    datos exceden sistemáticamente esa cota en $\sim8\%$
    ($N_c^{\rm libre}\approx0.8$--$1.1\,N$; ver Sección~\ref{sec:sweep_N}).
    Usando $\alpha_{T_B}$ solo, $N_c\propto N^{0.20}$ no converge a una
    interpretación simple.

    \item \textbf{$N=13$ está en un régimen distinto.}
    No alcanza el plateau de equipartición en $v\leq2.0\,\sigma/\tau$:
    la rigidez de la capa FCC completa más pequeña impide la termalización
    interna rápida.

    \item \textbf{La escala $h_1\approx T_B/(Mg)$ es válida para el primer
    rebote (ratio$=0.96$) pero no como attractor.}
    La altura CM decae geométricamente mientras $T_B$ decrece más despacio,
    y la bola queda atrapada sobre la superficie ($h_n\lesssim1\,\sigma$)
    tras $\sim23$ rebotes ($t\approx660\,\tau$ para $h_0=28\,\sigma$, $N=79$).
    La posterior filtración de átomos a través de la red
    ($t\approx713\,\tau$) es un artefacto del potencial repulsivo puro,
    no parte de la física del atrapamiento.

    \item \textbf{Con adhesión ($c_{ps}=1$), la bola pega en todos los casos
    estudiados ($v$ hasta $2.0\,\sigma/\tau$).}
    $E_{\rm adh}\propto N^{0.73}$, $h_0^*(N)\propto N^{-0.27}$.
    El umbral extrapolado $h_0^*\sim260$--$440\,\sigma$ es una cota de orden
    de magnitud (probablemente inferior; ver Sección~\ref{sec:adhesion});
    la escala favorable a cúmulos grandes tiene implicaciones para el
    \emph{soft landing}.

\end{enumerate}

Como trabajo futuro, la extensión a temperatura inicial $\Tstar>0$ y a
múltiples orientaciones del cúmulo (3--5 réplicas por punto) permitirá
asignar barras de error a las fracciones $\alpha_i$ ---y decidir si la no
monotonía de $e_1$ con $h_0$ (Tabla~\ref{tab:particion4}) es física o
varianza de condición inicial---, cuantificar la dependencia estadística de
$n_{\rm adh}$ (número de rebotes hasta adhesión) y el papel de la fusión
interna en la cascada. La medición del momento angular post-rebote
(\texttt{compute angmom}) separará la contribución rotacional de
$\alpha_{\rm bola}$, y el barrido de las capas completas intermedias
$N=19$ y $43$ poblará la transición molecular$\to$bulk entre
$N=13$ y $N=55$.

% =============================================================================
\section*{Agradecimientos}
% =============================================================================

Los autores agradecen al programa de Magíster en Simulación Computacional de la
PUCV por el apoyo académico. Las simulaciones se realizaron con
\lammps~7~Feb~2024~\cite{Thompson2022}. El código fuente está disponible en
\url{https://github.com/fwasaff/inelastic-rebound-md}.

% =============================================================================
\bibliographystyle{unsrt}
\begin{thebibliography}{99}

\bibitem{Gutierrez2012}
G. Gutiérrez, S. Davis, C. Loyola, J. Peralta, F. González, Y. Navarrete
y F. González-Wasaff,
\textit{Inelastic Collisions and Hypervelocity Impacts at Nanoscopic Level:
A Molecular Dynamics Study},
en: \textit{Molecular Dynamics -- Theoretical Developments and Applications
in Nanotechnology and Energy}, cap.~12, pp.~229--248, InTech (2012).

\bibitem{Haff1983}
P.~K. Haff,
\textit{Grain flow as a fluid-mechanical phenomenon},
J. Fluid Mech. \textbf{134}, 401 (1983).

\bibitem{Dahneke1975}
B. Dahneke,
\textit{Further measurements of the bouncing of small latex spheres},
J. Colloid Interface Sci. \textbf{51}, 58 (1975).

\bibitem{Awasthi2007}
A. Awasthi, S.~C. Hendy, P. Zoontjens y S.~A. Brown,
\textit{Molecular dynamics simulations of reflection and adhesion behavior in
Lennard-Jones cluster deposition},
Phys. Rev. B \textbf{76}, 115437 (2007).

\bibitem{Awasthi2010}
A. Awasthi, S.~A. Brown y S.~C. Hendy,
\textit{Oblique impacts of Lennard-Jones clusters},
Phys. Rev. B \textbf{82}, 073406 (2010).

\bibitem{Northby1987}
J.~A. Northby,
\textit{Structure and binding of Lennard-Jones clusters: $13 \le N \le 147$},
J. Chem. Phys. \textbf{87}, 6166 (1987).

\bibitem{Wales1997}
D.~J. Wales y J.~P.~K. Doye,
\textit{Global optimization by basin-hopping and the lowest energy structures
of Lennard-Jones clusters containing up to 110 atoms},
J. Phys. Chem. A \textbf{101}, 5111 (1997).

\bibitem{Haberland1994}
H. Haberland, M. Mall, M. Moseler, Y. Qiang, T. Reiners y Y. Thurner,
\textit{Filling of micron-sized contact holes with copper by energetic cluster
impact},
J. Vac. Sci. Technol. A \textbf{12}, 2925 (1994).

\bibitem{Thompson2022}
A.~P. Thompson \textit{et al.},
\textit{LAMMPS -- A flexible simulation tool for particle-based materials
modeling},
Comput. Phys. Commun. \textbf{271}, 108171 (2022).

\end{thebibliography}

\end{document}
